% !TEX root = main.tex
\section{A Dynamical View on ABC conjecture}\label{sec:a-dynamical-view-on-abc-conjecture}

In \cref{theorem:inductive-formula}, we can find the exponent of $\alpha$ is always $n$ wihle exponent of $\beta$ is variable with reduce sequence. Given a sequence, we also have its Cauchy production
\begin{definition}[Cauchy production]
	Given a sequences $\{a_k\}\subset \mathbb{Z}^{*}$. Define $\mathcal{A}_{n}^{m} = \prod_{k=m}^{n-1} a_k$ as its Cauchy production.
\end{definition}
\begin{example}
	Given $\alpha \in \mathbb{N}_{\ge2}$ are coprime. Given two sequences $\{a_k\}_{k=1}^{n} \subset\mathbb{N}$, then we have a Cauchy production $\mathcal{A}_{n}^{m} = \alpha^{\sum_{k=m}^{n-1}a_k}$.
\end{example}
\begin{theorem}\label{thm:power-diff-system-exists}
	Given two sequences $\{a_k\}_{k=1}^{n}, \{b_k\}_{k=1}^{n}\subset\mathbb{N}_{\ge1}$, we have their production $\mathcal{A}$ and $\mathcal{B}$. Given $r\in \mathbb{Z}^{*}$, then $\exists \{\gamma_k\}_{k=1}^{n}\subset\mathbb{Z}^{*}$ and $b \in\mathbb{Z}^{*}$, such that
	\[ s\cdot\mathcal{B}_n - r\cdot\mathcal{A}_n = \sum_{k=0}^{n-1}\mathcal{B}_k\mathcal{A}_{n-1}^{k}\gamma_k. \]
	We call these systems as general power-difference system.
\end{theorem}
\begin{proof}
	Similar to \cref{theorem:inductive-formula}. Just let $\gamma_k$ and $a_k$ satisfy
	\[ r_{k+1}\cdot b_{k} = r_{k}\cdot a_k + \gamma_k, \]
	and let $s = r_n$, the theorem holds.
\end{proof}

Also, we denote $ \mathfrak{U}_{n} = \sum_{k=0}^{n-1}\mathcal{B}_k\mathcal{A}_{n-1}^{k}\gamma_k$. We have $\mathfrak{U}\in\partial_{*}(\mathcal{B}_{*}, \mathcal{A}_{*})$ follows conclusions in \cref{sec:a-kind-of-derivation-on-groupoid} and its kernel is exactly $\gamma_{k}$. 
% \begin{remark}
% 	Notice that the system in \cref{thm:power-diff-system-exists} is built from reduce sequence, thus we have no result like $\beta \nmid b$ and $\beta\nmid \mathfrak{U}_n$. And we can't ensure $r_k > 0$ or $\gamma_k > 0$ any more.
% % \end{remark}

Not only these. From \cref{exa:fundamental-theorem-of-arithmetic-to-cauchy-homo}, alternatively, we have
\begin{theorem}\label{thm:arithmetic-decomposition-to-dynamical-system}
	For given $\mathcal{P}, \mathcal{Q}, \mathcal{S} \in \mathbb{N}$ that $\mathcal{Q} - \mathcal{P} = \mathcal{S}$.\newline
	When $\mathcal{P} = \prod_{k=0}^{n-1}p_k$ and $\mathcal{Q} = \prod_{k=0}^{n-1}q_k$, there are $\{r_k\}, \{\gamma_k\} \subset\mathbb{Z}^{*}$ of length $n$ such that $\mathcal{S}\cdot r_0 = \sum_{k=0}^{n-1}\mathcal{Q}_k\mathcal{P}_{n-1}^{k}\gamma_k$. In this condition, we have
	\begin{equation}\label{equ:equation-system-of-arithmetic-to-dynamical-system}
		\left\{\begin{aligned}
		&q_{0}\cdot r_{1} - p_{0} \cdot r_{0} = \gamma_{0},\\
		&q_{1}\cdot r_{2} - p_{1} \cdot r_{1} = \gamma_{1},\\
		% &q_{2}\cdot r_{3} - p_{2} \cdot r_{2} = \gamma_{2},\\
		&\cdots \\
		&q_{n-1}\cdot r_{0} - p_{n-1} \cdot r_{n-1} = \gamma_{n-1}.
	\end{aligned}\right.
	\end{equation}
\end{theorem}
\begin{proof}
	$\Longrightarrow$: We can always partition $\mathcal{P}$, $\mathcal{Q}$ at first, thus we suppose the sequences $\{p_k\}, \{q_k\}$ are already exists. From \cref{thm:power-diff-system-exists}, we have a such  $\{r_k\}, \{\gamma_k\} \subset\mathbb{Z}^{*}$. If for given $r_0 > 0$, there's $r_n \ne r_0$, then let $r_n \cdot \mathcal{Q} - r_0 \cdot \mathcal{P} = \mathfrak{U}$. With this we have $(r_n - r_0)\cdot \mathcal{Q} = \mathfrak{U} - r_0\cdot \mathcal{S}$. Now let
	\[ \gamma'_n = \gamma_n - (r_n - r_0)\cdot \mathcal{Q} / \mathcal{Q}_{n-1} = \gamma_n -  (r_n - r_0)\cdot q_{n-1}. \]
	From \cref{thm:power-diff-system-exists}, we can derive $r'_n$ with $\gamma'_n$ by
	\begin{align*}
		r'_n\cdot q_{n-1} &= a\cdot p_{n-1} + \gamma'_n \\
		&= a\cdot p_{n-1} + \gamma_n - (r_n - r_0)\cdot q_{n-1} \\
		&= r_n \cdot q_{n-1} - (r_n - r_0)\cdot q_{n-1} \\
		&= r_0 \cdot q_{n-1},
	\end{align*}
	which means we found a new $\{\gamma'_k\}$ and $\{r_k\}$ such that $r'_n = r_0$.
	\newline
	$\Longleftarrow$: Trivial from \cref{thm:derivation-periodic-rule}.
\end{proof}
\begin{remark}
	we can always suppose $\gcd(\gamma_0, \gamma_1, \cdots, \gamma_{n-1}) = 1$. And suppose $\forall k < n$, $p_k$, $q_k$ not both $1$.
\end{remark}
Actually, if $a_{j} = b_{j} = 0$. Then from $q_{j}\cdot r_{j+1} - p_{j} \cdot r_{j} = r_{j+1} - r_{j} = \gamma_{j}$ we have $q_{j+1}\cdot r_{j+2} - p_{j+1} \cdot r_{i} = \gamma_{j+1} + p_{j+1}\gamma_{i}$. Then let
\begin{align*}
	u_k &= \left\{\begin{aligned}
		&a_{k},  &&k < j, \\
		&a_{k+1}, && j < k < n-1,
	\end{aligned}\right.\\
	v_k &= \left\{\begin{aligned}
		&b_{k},  &&k < j, \\
		&b_{k+1}, && j < k < n-1,
	\end{aligned}\right.\\
	\eta_k &= \left\{\begin{aligned}
		&\gamma_{k},  &&k < j, \\
		&\gamma_{j+1} + p_{j+1}\gamma_{j}, && k = j+1, \\
		&\gamma_{k+1}, && j+1 < k < n-1.
	\end{aligned}\right.
\end{align*}
We build a same system of equations for $\{u_k\}, \{v_k\}, \{r_k\}, \{\eta_k\}$ like \cref{equ:equation-system-of-arithmetic-to-dynamical-system}, and $u_j, v_j$ are not both $0$. Repeat this operation, we can replace all $p_{j} = q_{j} = 1$, that's why we can always suppose $p_k$, $q_k$ not both $1$. From this operation, we can also find that if \cref{equ:equation-system-of-arithmetic-to-dynamical-system} exists for a given $N$, then similar systems of equations also exist for all $n$ that $0 < n \le N$.

\cref{thm:arithmetic-decomposition-to-dynamical-system} reveals the connection between $\mathcal{S}$ and the factors of $\mathcal{P}$ and $\mathcal{Q}$ in the equation $\mathcal{P} + \mathcal{S} = \mathcal{Q}$; this bears a close relationship to the ABC conjecture. Actually, with this theorem, we have
\begin{conjecture}[rewitten ABC conjecture]
	$\forall \varepsilon \ge 0$, there exist only finite triple $(\mathcal{P}, \mathcal{Q})$ such that
	\[ \rad\left(\frac{1}{r_0}\sum_{k=0}^{n-1}\mathcal{Q}_k\mathcal{P}_{n-1}^{k}\gamma_k\right) < \frac{\mathcal{Q}^{1-\varepsilon}}{\rad(\mathcal{P}\mathcal{Q})} = \frac{\mathcal{Q}^{\sharp-\varepsilon}}{\rad\mathcal{P}}. \]
	Here $\mathcal{Q}^{\sharp} = \mathcal{Q} / \rad\mathcal{Q} = \partial\mathcal{Q}/\partial\!\rad\mathcal{Q}$.
\end{conjecture}

% Suppose $p_k$ and $q_k$ are both prime numbers or $1$.

\begin{corollary}\label{cor:power-diff-to-dynamical-system}
	Suppose $\alpha, \beta \in \mathbb{N}_{\ge2}$ are coprime. For given $C \in \mathbb{Z}$, the equation $\beta^{B} - \alpha^{A} = C$ possesses an integer solution $A, B$ if and only if $\exists n \in \mathbb{N}$ and sequences $\{a_k\}_{k=1}^{n}, \{b_k\}_{k=1}^{n} \subset \mathbb{N}$ and $\{r_k\}_{k=1}^{n}, \{\gamma_k\}_{k=1}^{n} \subset\mathbb{Z}^{*}$, such that whenever $A = \sum_{k=0}^{n-1}a_k$ and $B = \sum_{k=0}^{n-1}b_k$, there is $C \cdot r_0 = \sum_{k=0}^{n-1}\beta^{\mathcal{B}_k}\alpha^{\mathcal{A}_{n-1}^{k}}\gamma_k$, which also equivalents to the sequences $\{a_k\}, \{b_k\}, \{r_k\}, \{\gamma_k\}$ satisfying the following system of equations:
	\begin{equation}\label{equ:equation-system-of-power-diff-to-dynamical-system}
		\left\{\begin{aligned}
		&\beta^{b_{0}}\cdot r_{1} - \alpha^{a_{0}} \cdot r_{0} = \gamma_{0},\\
		&\beta^{b_{1}}\cdot r_{2} - \alpha^{a_{1}} \cdot r_{1} = \gamma_{1},\\
		% &\beta^{\mathcal{B}_{3}^{2}}\cdot r_{3} - \alpha^{\mathcal{A}_{n-1}^{2}} \cdot r_{2} = \gamma_{2},\\
		&\cdots \\
		&\beta^{b_{n-1}}\cdot r_{0} - \alpha^{a_{n-1}} \cdot r_{n-1} = \gamma_{n-1}.
	\end{aligned}\right.
	\end{equation}
\end{corollary}
\cref{cor:power-diff-to-dynamical-system} posits a perspective: given a constant $C$, we can derive solutions $A$ and $B$ by searching for partitions of the exponents.
Indeed, as established by the analysis in the preceding sections, every equation of the form presented in \cref{cor:power-diff-to-dynamical-system} corresponds to a specific "top branch" of a certain $t$-periodic power-difference dynamical system.
We can employ the methods outlined in \cref{sec:distribution-of-source-value-set} to enumerate the number of such top branches contained within this system. These top branches encompass all the cycles within the system—cycles that correspond precisely to the solutions of the aforementioned power-difference equation.
Since as defined in \cref{sec:a-kind-of-derivation-on-groupoid}, the specific $\gamma_k$ values we seek constitute a discrete, exact differential form; this fact avoids Conway's limitation and allows us to draw inspiration from the principles of homology theory to analyze the number-theoretic properties that these $\gamma_k$ values must satisfy.


\cref{thm:arithmetic-decomposition-to-dynamical-system} prompt us to investigate $\gamma_k$, the exact differential forms associated with such discrete dynamical systems. Naturally, these quantities are interrelated with the $r_k$, and thus they should be considered as well.

% TODO.

% \newpage
% Given $f(x) = k\prod_{i=0}^{m}(x - r_i)^{n_i}$. We have $f' = k\sum_{i=1}^{m}\frac{n_i f}{x - r_i}$.
% Define $\rad f = \prod_{i=1}^{m}(x - r_i)$. Since we have
% \[ \frac{f'}{f} = \sum_{i=1}^{m}\frac{n_i}{x - r_i} = \frac{1}{\rad f}\sum_{i=1}^{m}\frac{n_i\rad f}{x - r_i} \]
% Denote $\rad' f = \sum_{i=1}^{m}\frac{n_i\rad f}{x - r_i}$. Then we have
% \[ f'\cdot \rad f = f\cdot \rad'f. \]
% Notice that $\rad' f \ne (\rad f)'$. Actually, $(\rad f)' = \sum_{i=1}^{m}\frac{\rad f}{x - r_i}$ because the multiplicity of each root of $\rad f$ is $1$.

% To prove Mason-Stothers theorem, we define $W(f, g) = g'f - f'g$ is the Wronskian determinant. Consider $W(f, g)\cdot\rad(fg)$, we have 
% \begin{align*}
% 	W(f, g)\cdot\rad(fg) &= (g'f - f'g)\cdot\rad(fg)\\
% 	&= g'f\cdot\rad g\cdot\rad f - gf'\cdot\rad f\cdot\rad g \\
% 	&= g\rad'g\cdot f\rad f - f\rad'f\cdot g\rad g \\
% 	&= \left[g\rad'g\cdot f\rad f - f\rad'f\cdot g\rad g\right] \\
% 	&= \left[\rad'g\cdot \rad f - \rad'f\cdot \rad g\right]\cdot fg.
% \end{align*}
% Which means $fg \mid W(f, g)\rad(fg)$.
% For $f + g = h$ that $\gcd(f, g) = 1$, we have $W(f, g) = W(f, h) = W(h, g)$. Thus we have $fh \mid W(f, h)\rad(fh) = W(f, g)\rad(fh)$.
% Sum up, we have $fgh \mid W(f, g)\rad(fgh)$. Thus $\deg(fgh) \le \deg\left[W(f, g)\rad(fgh)\right]$, which equivalents to
% \[ \deg f + \deg g + \deg h \le \deg W(f, g) + \deg \rad(fgh). \]
% Since $\deg W(f, g) = \deg (g'f - f'g) \le \deg f + \deg g - 1$, we have
% \[ \deg h \le \deg \rad(fgh) - 1. \]
% Or more likely, we have
% \begin{theorem}[Mason-Stothers]
% 	$\max\{ \deg f, \deg g, \deg h \} \le \deg \rad(fgh) - 1$.
% \end{theorem}
% Notice that for a polynomial $f$, the $\deg f$ represents the number of its roots in an algebraically closed field, counting multiplicities. Thus, this theorem provides a limitation on the multiplicities of the polynomial's roots.

\begin{definition}\label{def:der-rad-rings}
	Suppose $\mathcal{C} = R(\mathbb{K})/\mathbb{K}$ is a category of integral domain modulo the algebrical closed $\mathbb{K}$ characteristic $0$. We require $\rad \in \mathrm{Mor}(\mathcal{C})$ and $\partial, \overline{\partial}: \mathcal{C} \to \mathcal{C}$, they satisfy:
	\begin{itemize}
		\item[1.] $\rad(f^{2}) = \rad{f}$ and $(\rad f) \mid f$. We denote $f^{\sharp} = f/\rad f$.
		\item[2.] $\rad (f\cdot g)\cdot\rad\gcd(f, g) = \rad(f)\cdot\rad(g)$.
		\item[3.] $\gcd(\rad f, \rad g) = \rad\gcd(f, g)$.
		\item[4.] $\partial a = 0 \Longleftrightarrow a = 1$.
		\item[5.] $\partial(f\cdot g) = f\cdot\partial{g} + g\cdot\partial{f}$ and $\overline{\partial}(f \cdot g) = f\cdot\overline{\partial}{g} - g\cdot\overline{\partial}{f}$.
		\item[6.] $\rad \overset{\partial}{\mapsto} \partial\!\rad$ and $\rad \overset{\overline{\partial}}{\mapsto} \overline{\partial}\!\rad$ such that
		\begin{align*}
			\partial f \cdot \rad f &= f \cdot \partial\!\rad f, \\
			\overline{\partial} f \cdot \rad f &= f \cdot \overline{\partial}\!\rad f.
		\end{align*}
	\end{itemize}
\end{definition}
\begin{remark}
	In $\partial(f\cdot g)$, the dot ($\cdot$) represent the ring multiplication. However, in $\overline{\partial}(f\cdot g)$, the dot ($\cdot$) is no longer ring multiplication but merely a separator to consistency in form with $\partial$. Even though, we still using $\partial(fg)$ and $\overline{\partial}(fg)$, they still represent $\partial(f\cdot g)$ and $\overline{\partial}(f\cdot g)$. Actually, $\overline{\partial}$ is the abstraction of Wronskian determinant.
\end{remark}
With this we have
\begin{proposition}
	$\overline{\partial}(f\cdot\overline{\partial}(g\cdot h)) + \overline{\partial}(g\cdot\overline{\partial}(h\cdot f)) + \overline{\partial}(h\cdot\overline{\partial}(f\cdot g)) = 0$.
\end{proposition}
Denote $\widehat{n} = (-1)^{\mathfrak{p}}(n\fmod 2)$, here $\mathfrak{p}(i,j)$ is the character of path that $i + j = n$. With this we have
\begin{proposition}
	$\partial(f^{n}) = nf^{n-1}\partial f$ and $\overline{\partial}(f^{n}) = \widehat{n}f^{n-1}\overline{\partial}f$.
\end{proposition}
% \begin{proof}
	
% \end{proof}
\begin{proposition}\label{prop:operator-are-separable}
	Suppose $f = \prod_{i=1}^{m}f_{i}^{n_i}$ that $f_i\ne 1, 0$.  We have
	\begin{align*}
		\partial f &= \sum_{i=1}^{m}\frac{f}{f_i}n_i\partial{f_i}, & \partial\!\rad f &= \sum_{i=1}^{m}\frac{\rad f}{f_i}n_i\partial\!\rad{f_i}.
	\end{align*}
	When right associative, $f = f_1^{n_1}\cdot (f_2^{n_2}\cdot (\cdots f_m^{n_m}))$, we have
	\begin{align*}
		\overline{\partial}f &= \sum_{i=1}^{m}\frac{f}{f_i}(-1)^{i-1}\widehat{n_i}\overline{\partial}{f_i}, & \overline{\partial}\!\rad f  &= \sum_{i=1}^{m}\frac{\rad f}{f_i}(-1)^{i-1}\widehat{n_i}\overline{\partial}\!\rad{f_i}.
	\end{align*}
	When left associative, $f = ((f_1^{n_1}\cdot f_2^{n_2}) \cdots) \cdot f_m^{n_m}$, we have
	\begin{align*}
		\overline{\partial}f &= \sum_{i=1}^{m}\frac{f}{f_i}(-1)^{m-i}\widehat{n_i}\overline{\partial}{f_i}, &\overline{\partial}\!\rad f  &= \sum_{i=1}^{m}\frac{\rad f}{f_i}(-1)^{m-i}\widehat{n_i}\overline{\partial}\!\rad{f_i}.
	\end{align*}
	Finally, we have $\rad f = \prod_{i=1}^{m}\rad f_{i}$.
\end{proposition}
Generally, we have $\partial\!\rad f \ne \partial(\rad f)$ and $\overline{\partial}\!\rad f \ne \overline{\partial}(\rad f)$. Not only this, since we have no $\overline{\partial}(f+g) = \overline{\partial} f + \overline{\partial} g$, we can't tell that if $\overline{\partial}(f\cdot g) = 0$
\begin{proposition}
	If $\gcd(f, g) = 1$, then
	\begin{align*}
		\partial\!\rad(fg) &= \rad f\cdot \partial\!\rad g + \rad g\cdot \partial\!\rad f, \\
		\overline{\partial}\!\rad(fg) &= \rad f\cdot \overline{\partial}\!\rad g - \rad g\cdot \overline{\partial}\!\rad f.
	\end{align*}
\end{proposition}

Since there is no $\overline{\partial}(f + g) = \overline{\partial}f + \overline{\partial}g$ in our axioms, we cannot obtain a conclusion analogous to the Mason-Stothers theorem. However, we still expect these deviations not to be too significant.
Therefore, we denote that
\begin{align*}
	\delta\partial(f + g) &= \partial(f + g) - \partial f - \partial g,\\
	\delta\overline{\partial}(f + g) &= \overline{\partial}(f + g) - \overline{\partial}f - \overline{\partial}g.
\end{align*}
Of course, here ($+$) also means separator instead of ring addition. That means even $s + t  = f + g$, we may have no $\delta\overline{\partial}(s + t) = \delta\overline{\partial}(f + g)$.
Similar with Mason-Stothers theorem, we have
\begin{proposition}\label{prop:fgh-mid-lcm-rad}
	If $f + g = h$ and $\gcd(f, g) = 1$, then we have
	\[ fgh\ \left\vert\ \lcm\left\{\overline{\partial}(fg), \overline{\partial}(fg) + f\delta\overline{\partial}(f + g)\right\}\cdot\rad(fgh). \right. \]
	And thus
	\[ f^{\sharp}g^{\sharp}h^{\sharp}\ \left\vert\ \lcm\left\{\overline{\partial}(fg), \overline{\partial}(fg) + f\delta\overline{\partial}(f + g)\right\} \right. \]
\end{proposition}
\begin{proof}
	At first, we have
	\begin{align*}
		\overline{\partial}(fg) \cdot \rad(fg) &= fg \cdot \overline{\partial}\!\rad(fg), \\
		\overline{\partial}(fh) \cdot \rad(fh) &= fh \cdot \overline{\partial}\!\rad(fh).
	\end{align*}
	Thus $fgh \left\vert\ \lcm(\overline{\partial}(fg), \overline{\partial}(fh))\cdot\rad(fgh)\right.$. And we have
	\begin{align*}
		\overline{\partial}(fh) &= f\overline{\partial}(h) - h\overline{\partial}(f) \\
		&= f\delta\overline{\partial}(f + g) + f\overline{\partial}(f) + f\overline{\partial}(g) - f\overline{\partial}(f) - g\overline{\partial}(f) \\
		&= \overline{\partial}(fg) + f\delta\overline{\partial}(f + g).
	\end{align*}
	With this we proved this proposition.
\end{proof}
This proposition indicates that $\delta\overline{\partial}(f + g)$ cannot be chaotically random; it is strictly controlled by $f^{\sharp}g^{\sharp}h^{\sharp}$.
In practice, we want $\vert\delta\overline{\partial}(f + g)\vert$ to be as small as possible; this ensures that our inequality is as tight as possible.
From \cref{prop:operator-are-separable}, we can only consider derivations on those inseparable elements, since derivations on other elements can be automatically generated by those on inseparable elements, in accordance with our axioms.

With this proposition, we denote
\[ \mathcal{L}(f, g) = \lcm\left\{\overline{\partial}(fg), \overline{\partial}(fg) + f\delta\overline{\partial}(f + g)\right\}. \]
Additionally, denote
\begin{align*}
	\mathcal{M}(f, g) &= \frac{\mathcal{L}(f,g)}{f^{\sharp}g^{\sharp}h^{\sharp}} \in \mathrm{Obj}(\mathcal{C}), &\mathcal{K}(f, g) &= \frac{\mathcal{L}(f,g)}{f\cdot g} = \frac{h^{\sharp}\cdot\mathcal{M}(f,g)}{\rad(f\cdot g)}.
\end{align*}
Here, we do not require $\mathcal{K}(f, g) \in \mathrm{Obj}(\mathcal{C})$, but it should be the natural localization domain of the rings.
For polynomials, the $\mathcal{K}(f, g) = \frac{f'}{f} - \frac{g'}{g}$. However, for natural numbers, we don't know yet what the $\mathcal{K}(f, g)$ and $\mathcal{M}(a, b)$ represent for.

We have to say, here $\overline{\partial}(h)$ is not necessarily equal to $f\overline{\partial}g - g\overline{\partial}f$. In fact, $\overline{\partial}(h)$ should be viewed as a set rather than a single value. While we can always guarantee that $f\overline{\partial}g - g\overline{\partial}f \subset \overline{\partial} h$, this does not imply equality between the two. Unlike the case of integers or certain integral domains—where $h$ possesses a unique factorization—with the operator $\overline{\partial}$, both the granularity of the factorization and the order of the resulting factors affect the final outcome. However, this does not imply that the concept is ill-defined or that Proposition 1 is absurd. Indeed, when we ensure $f + g = h$, the variation in factorization results affects only $\mathcal{M}(f, g)$; this actually suggests that Proposition 1 not only holds but also reflects a profound property of factorization in integral domains that we do not yet fully understand.

Since the classical arithmetic derivation that $\partial(a) = a\sum_{p\in\mathbb{P}}^{p\mid a}\frac{\val_{p}(a)}{p}$ with initial $\overline{\partial}(1) = \partial(1) = 0$ and $\overline{\partial}(p) = \partial(p) = 1$ exactly follows the axioms in \cref{def:der-rad-rings}.
Equivalents to the ABC conjecture, we have
\begin{conjecture}
	$\forall \varepsilon > 0$, $\exists K_{\varepsilon} > 0$, such that $\forall a, b, c\in\mathbb{N}_{\ge 1}$ that $\gcd(a, b, c) = 1$ and $a + b = c$, there is $\frac{\mathcal{K}(a, b)}{\mathcal{M}(a, b)} = \frac{c}{\rad(abc)} \le K_{\varepsilon}\cdot \rad(abc)^{\varepsilon}$.
\end{conjecture}
